\subsection{A complex network from orthogonal pair-star sums}
\label{sec:pair-star-complex}
The following construction changes the complex scalar network and its binary
labels.  The bit network used for address movement is unchanged.  Let $h>6$
be even, use ground coordinates $\{1,\ldots,h\}$, let
$\mathcal T=\binom{\{1,\ldots,h\}}3$, and put
\[
 v=\binom h3,\qquad N=v^3,\qquad m=h^3,\qquad
 F=\mathbb F_2^h,\qquad \mathcal F=F^{\otimes3}.
\]
The form on $F$ is the ordinary binary dot product.  For a triple $T$,
write $t_T$ for its indicator.  Thus $t_T$ has binary norm one, and
$t_S\cdot t_T=0$ whenever $|S\cap T|$ is even.  All scalar arithmetic
below takes place in $\mathbb Z[i,1/2]$; the binary field is used only
for labels and address frames.
The two data banks $X_a,Y_a$ have indices
$a=(a_1,a_2,a_3)\in\mathcal T^3$.  Write
$u_a=t_{a_1}\otimes t_{a_2}\otimes t_{a_3}$ and
$U_a=\langle u_a\rangle$.  The choice of one-based ground coordinates
is only a relabeling of the original interface.

\paragraph{Partitioning the side corrections.}
Give the ground set its natural order.  Write each source triple uniquely as
$T=\{a,b,k\}$ with $a<b<k$.  For each fixed pair $(a,b)$, partition the
ordered list of possible $k>b$ into consecutive batches, each of size at
most $h-6$.  These batches partition the source triples.  Fix one batch
$K$ and a target triple $S$, and put $r=|S\cap\{a,b\}|$.  The contribution
of this batch to the side injection is
\[
 \begin{cases}
 \displaystyle\frac12\sum_{k\in K\setminus S}x_{\{a,b,k\}},&r=0,\\[5pt]
 \displaystyle-\frac12\sum_{k\in K\cap S}x_{\{a,b,k\}},&r=1,\\[5pt]
 \displaystyle-\frac12\sum_{k\in K\setminus S}x_{\{a,b,k\}},&r=2.
 \end{cases}                                                    \tag{PS1}
\]
An empty sum is omitted.  For $r=0$, only intersection-zero sources
survive; for $r=1$, only intersection-two sources survive; for $r=2$,
the intersection-two sources survive and the possible diagonal source
$T=S$ is excluded.  Summing (PS1) over the disjoint source batches therefore
gives exactly
\[
 -\sum_{\substack{T\ne S\\ |S\cap T|\ {\rm even}}}
                  \frac{|S\cap T|-1}{2}\,x_T.                    \tag{PS2}
\]
The signs in (PS1) belong to the output injections.  All intermediate
nodes compute sums with coefficient $+1$.

\paragraph{A specified cancellation-free sum graph.}
For each batch, represent a nonempty source subset by a bit mask $M$, with
the bits ordered by increasing $k$.  A singleton mask is an input node.
For a mask with at least two set bits, let $\ell$ be the lowest set-bit
index and let $u$ be one more than the highest set-bit index.  Set
$j=\lfloor(\ell+u)/2\rfloor$ and split $M$ into its bits below $j$ and
its bits at least $j$.  Both pieces are nonempty and disjoint.  Compute
their sums recursively and add the results.  Intern every mask within
this batch, so a previously computed formal sum is reused.  Process
batches in the order of their anchor pair and then their first $k$, and
process target triples in lexicographic order.  Request precisely the
nonempty supports in (PS1), keeping one output use per target and batch,
even when two uses request the same sum.  This completely specifies the
finite graph; its children are defined using the occupied-mask endpoints,
not the endpoints of the entire batch.

Every addition has disjoint input supports.  Consequently a source has
coefficient one in a node exactly when it belongs to that node's mask;
there is no cancellation and no repeated contribution.  Let $c_{\rm add}$
be the number of distinct addition nodes and let $q_{\rm out}$ count
output uses, including multiplicities.  The same graph is used for every
invocation of a stage.  Its input triples are partitioned between batches,
so no node combines different batches.

\paragraph{Reversible roles and signed restoration.}
Treat each use by an addition and each designated output use as an outgoing
edge.  At an input node allocate one role for every outgoing edge, select
one as a pivot, and add the pivot into the others.  At an addition node
reuse its first incoming role as the pivot for one outgoing edge; allocate
a new role for every other outgoing edge.  Add the second incoming value
to the pivot and then add the pivot into the new outgoing roles.  The
second incoming role is thereafter inactive during the forward graph.
Every node operation is invertible by reversing its additions.  Roles
are identified only along directed paths, so two simultaneously used
incoming roles remain distinct.

There are $2c_{\rm add}+q_{\rm out}$ outgoing uses and precisely one
role identification per addition.  Thus the side computation uses
\[
 R_{\rm side}=c_{\rm add}+q_{\rm out}                            \tag{PS3}
\]
roles.  Write $\mathcal L$ for the resulting invertible mixer on these
roles, $\mathcal V$ for copying the physical sources into the input
pivots, and $\mathcal J$ for injecting the designated output slots into
their targets with the coefficients in (PS1).  On a zero side vector,
$\mathcal J\mathcal L\mathcal V$ is (PS2).

Let $\mathcal G,\mathcal R$ denote the central gather and scatter defined
below.  An operation written $+\mathcal J$ or $-\mathcal J$ means adding
or subtracting the indicated injection into the targets, leaving all its
controls unchanged; the same convention applies to $\mathcal V$,
$\mathcal G$ and $\mathcal R$.  The forward chronological schedule is
\[
 \mathcal L,\ -\mathcal J,\ \mathcal L^{-1},\ -\mathcal R,\
 +\mathcal V,\ +\mathcal G,\ +\mathcal R,\
 \mathcal L,\ +\mathcal J,\ \mathcal L^{-1},\
 -\mathcal G,\ -\mathcal V.                                   \tag{PS4}
\]
For arbitrary initial side vector $z$ and central vector $c_0$, the net
target change is
\[
 -\mathcal J\mathcal Lz-\mathcal Rc_0
 +\mathcal R(c_0+\mathcal Gx)
 +\mathcal J\mathcal L(z+\mathcal Vx)
 =(\mathcal J\mathcal L\mathcal V+\mathcal R\mathcal G)x.
\]
The final two operations restore the centers and the side vector exactly.
The minus signs in (PS4) are essential over the complex scalar ring.

\paragraph{Only $h$ central coordinates are needed.}
Use center roles $C_1,\ldots,C_{h-1},C_\ast$.  Gather by adding
\[
 (\mathcal Gx)_i=\sum_{T\ni i}x_T\quad(i<h),\qquad
 (\mathcal Gx)_\ast=\sum_Tx_T.
\]
Define the scatter at target $S$ by
\[
 (\mathcal Rc)_S=
 \begin{cases}
 \displaystyle\frac12\left(\sum_{i\in S}c_i-c_\ast\right),
                                    &h\notin S,\\[6pt]
 \displaystyle c_\ast-\frac12
                  \sum_{\substack{i<h\\i\notin S}}c_i,
                                    &h\in S.
 \end{cases}                                                     \tag{PS5}
\]
The omitted gather coordinate equals
$3(\mathcal Gx)_\ast-\sum_{i<h}(\mathcal Gx)_i$.  Substitution into
the original scatter proves
\[
 (\mathcal R\mathcal G)_{S,T}=\frac{|S\cap T|-1}{2}.
\]
This identity is asserted for gather increments, not for arbitrary initial
center values.  The latter cancel in (PS4), so no initial consistency
condition is imposed on the centers.  Combining this identity with (PS2)
gives $\mathcal J\mathcal L\mathcal V+\mathcal R\mathcal G=I$.
Thus every forward invocation is $y\gets y+x$ and restores arbitrary
scratch.  The inverse invocation uses the reversed inverse schedule.
The three stage orientations remain forward $X\to Y$, inverse $Y\to X$,
and forward $X\to Y$, and hence send $(X,Y)$ to $(-Y,X)$.

\paragraph{Common frames in the two orientations.}
Use the original stage decomposition
\[
 A=F^{\otimes(j-1)}=B\mathbin\perp P,
 \qquad P=\langle t_{a_1}\otimes\cdots\otimes t_{a_{j-1}}\rangle,
 \qquad Q=\langle t_{a_{j+1}}\otimes\cdots\otimes t_{a_3}\rangle.
\]
Empty tensor products are the one-dimensional field.  Both $P$ and $Q$
have norm-one generators.  Suppress $Q$ temporarily and write
\[
 D_0=B\otimes F,\quad D_1=A\otimes F,\quad
 D_U=(B\otimes F)\mathbin\perp(P\otimes U),
 \quad X_t=D_{\langle t\rangle},\quad Y_t=D_{t^\perp}.
\]
All these labels are tensored with $Q$ in $\mathcal F$.

At a graph node $z$, let $U_z$ be the span of the source indicators in
its mask.  Two distinct triples in its batch share exactly the two anchors,
so their binary indicators are orthogonal.  Each has norm one.  Thus
these indicators form an orthonormal basis of $U_z$.  Every $U_z$ is
nondegenerate, and along an edge $z\to w$ one has $U_z\subset U_w$.
At a designated output for target $S$, all contributing source indicators
are orthogonal to $t_S$, by (PS1); hence $U_z\subset t_S^\perp$.

The common frames for the twelve forward operations in (PS4) are, in order,
\[
 D_0,\ D_0,\ D_0,\ D_0,\ X_{t_T},\ D_1,\ D_0,\
 D_{U_z},\ Y_{t_S},\ D_1,\ D_1,\ D_1.                          \tag{PS6}
\]
The $X$ copy gates and $Y$ injection gates are grouped by their physical
data triple; the middle mixer is grouped by graph node.  At each such
gate the same displayed label is assigned to every incident role.
All early mixer gates use $D_0$ and all late cleanup gates use $D_1$.
An input or output node that otherwise gives an identity operation may
be retained as an identity gate to specify these incidences.

In the reversed invocation the middle computation is $\mathcal L^{-1}$.
Assign its node $z$ the complement frame $D_{U_z^\perp}$.  Reversing
an edge gives $U_w^\perp\subset U_z^\perp$.  At an original output
$\langle t_S\rangle\subset U_z^\perp$, and at an original input
$U_z^\perp=t_T^\perp$.  These are precisely the reversed physical
$X$ and $Y$ attachments.  In chronological order the reversed operations
and labels are
\[
\begin{array}{c|c|c@{\qquad}c|c|c}
\text{step}&\text{operation}&\text{label}
 &\text{step}&\text{operation}&\text{label}\\ \hline
1&+\mathcal V&D_0&7&-\mathcal G&D_0\\
2&+\mathcal G&D_0&8&-\mathcal V&Y_{t_T}\\
3&\mathcal L&D_0&9&+\mathcal R&D_1\\
4&-\mathcal J&X_{t_S}&10&\mathcal L&D_1\\
5&\mathcal L^{-1}&D_{U_z^\perp}&11&+\mathcal J&D_1\\
6&-\mathcal R&D_1&12&\mathcal L^{-1}&D_1.
\end{array}                                                     \tag{PS7}
\]
The logical source in (PS7) is the physical $Y$ bank.  Pointwise scalar
coefficients, including the signs and halves, commute with the common
array frame exactly as in the original frame identity.

\paragraph{Binary residuals, including entering and retired roles.}
The following argument checks the additional condition needed by the
complex phase transfer: every nonzero residual must have an orthonormal
basis, not merely be nondegenerate.  The source support of a node occupies
at most $(h-6)+2=h-4$ ground coordinates.  Adjoining one target triple
occupies at most $h-1$ coordinates.  There is therefore a coordinate unit
outside either union.  Such a unit belongs to $U_z^\perp$, and in the
second case it also belongs to $t_S^\perp$; its norm is one.

All the following residuals are nondegenerate.  For nested nondegenerate
spaces this follows by orthogonally splitting the larger space into the
smaller one and its residual.  Suppressing the common factor $P$, the
new forward residual classes are
\[
\begin{array}{c|c|l}
\text{transition}&\text{residual}&\text{norm-one witness if nonzero}\\ \hline
D_0\to D_U&U&\text{a source indicator}\\
D_U\to D_V\ (U\subset V)&V\cap U^\perp
                              &\text{an indicator in the support difference}\\
D_U\to Y_{t_S}&t_S^\perp\cap U^\perp
                              &\text{a unit outside both supports}\\
D_U\to D_1&U^\perp&\text{a unit outside the source support}.
\end{array}                                                     \tag{PS8}
\]
The second row includes an input line entering a later sum.  It has the
indicators in $V$'s support but not $U$'s support as an orthonormal basis.
The first row also covers a newly allocated outgoing role; the last row
covers a second input retired after an addition.  A direct $D_0\to D_1$
transition has residual $F$.

For the reversed computation the corresponding complete list is
\[
\begin{array}{c|c|l}
\text{transition}&\text{residual}&\text{norm-one witness if nonzero}\\ \hline
D_0\to D_{U^\perp}&U^\perp&\text{a unit outside the source support}\\
X_{t_S}\to D_{U^\perp}&t_S^\perp\cap U^\perp
                              &\text{a unit outside both supports}\\
D_{V^\perp}\to D_{U^\perp}\ (U\subset V)&V\cap U^\perp
                              &\text{an indicator in the support difference}\\
D_{U^\perp}\to D_1&U&\text{a source indicator}.
\end{array}                                                     \tag{PS9}
\]
Roles retired in the forward mixer become entering roles in its inverse,
and newly allocated forward roles become retired inverse roles.  Thus the
first and last rows of (PS9) are necessary; the two orientations do not
assume zero values on such roles.  At an original input the final attachment
to $Y_{t_T}$ is an equality of labels.  Tensoring a displayed witness with
the norm-one generators of $P$ and $Q$ gives a norm-one witness for the
actual residual in $\mathcal F$.

The remaining data and central residuals are those of the original tensor
stage table: zero spaces, copies of $P\otimes F$, copies of
$P\otimes t^\perp$, and copies of $B\otimes t^\perp$, all tensored
with $Q$.  Initial auxiliary edges have residual $B\otimes F\otimes Q$;
final auxiliary edges have residual $(A\otimes F)\otimes Q^\perp$.
When $B$ or $Q^\perp$ is nonzero it contains a coordinate unit outside
the support of its defining earlier or future tensor indicator.  Those
supports have size $3^j<h^j$.  A triple complement contains a coordinate
unit outside the triple.  These observations give norm-one witnesses for
every nonzero remaining tensor residual.  In particular all data boundaries
are the original ones.  The sole decreases are the $h$ central-role
returns from $D_1$ to $D_0$ in each invocation, each of dimension $h$.

A nondegenerate symmetric binary space containing a norm-one vector has an
orthonormal basis, by the elementary binary-space argument in the original
residual lemma.  Equations (PS8)--(PS9) and the preceding tensor witnesses
therefore verify the full orthonormal-residual hypothesis of the complex
phase interface.  No positivity property of the rational bit labels is
being substituted for this binary condition.

\paragraph{Reusing the first-stage auxiliary bank in stage three.}
Pair the $h$ ground coordinates consecutively and let $\pi$ send a triple
to its image under the permutation swapping the two members of every pair.
The intersection of a triple and its image has size zero or two, so
$t_A\cdot t_{\pi(A)}=0$.  The map $\pi$ is a bijection of triples.
Identify the entire auxiliary bank, including its $h$ centers, of the
stage-1 invocation with fixed triples $(A,B)$ with the auxiliary bank of
the stage-3 invocation with fixed triples $(B,\pi(A))$, retaining each
local role index.  This is a bijection of invocation banks.  The two
invocations occur at different stages and restore arbitrary scratch, so
no simultaneous values are identified.

Every auxiliary role now first has label $D_0\otimes Q$ and last has
label $D_1\otimes Q$.  The new interstage join is
\[
 E=F\otimes\langle t_A\rangle\otimes\langle t_B\rangle
 \ \subset\
 H=\langle t_B\otimes t_{\pi(A)}\rangle^\perp_{F\otimes F}
                                                 \otimes F.     \tag{PS10}
\]
Both spaces are nondegenerate, and the inclusion follows from
$t_A\perp t_{\pi(A)}$.  The residual $H\cap E^\perp$ is nondegenerate.
Choose a ground coordinate $j$ outside $A\cup\pi(A)$, and any coordinates
$r,k$.  The norm-one tensor $e_r\otimes e_j\otimes e_k$ belongs to
both $H$ and $E^\perp$.  Thus the new join residual also has an orthonormal
basis.  It introduces no decreasing edge.  It replaces the old edges
$E\to\mathcal F$ and $0\to H$ and reduces their total absolute
dimension change by exactly $m$ per removed role.

\paragraph{The coefficient-depth allowance still applies.}
Every elementary arithmetic operation uses only coefficients
$0,\pm1,\pm1/2$.  One application of $\mathcal L$ or its inverse
contains at most $2R_{\rm side}$ scalar additions or subtractions:
the outgoing-use count is $2c_{\rm add}+q_{\rm out}$, and the
input fanouts save one operation per input.  There are four mixer
applications in (PS4).  The two signed output injections cost at most
$4q_{\rm out}$ additions and halvings; the two source copies cost
at most $2v$ additions; the two gathers cost at most $8v$ additions;
and the two scatters cost at most $4vh$ additions and halvings.
Thus an invocation uses at most
\[
 12R_{\rm side}+4vh+10v
\]
elementary coefficient operations.  Reversing the schedule has the same
bound.  A conservative allowance for the complete three-stage scalar
circuit and its scalar endpoint corrections is therefore
\[
 D_{\rm scalar}
 =3v^2(12R_{\rm side}+4vh+10v)+8W.                              \tag{PS11a}
\]
This counts all operations in the finite scalar circuit, so in particular
it bounds every coefficient dependency path there.  For each inverse
recursive child, the two $Z^{\otimes f}$ corrections are each one sign
selected by address parity, and $(-i)^f$ costs at most a sign and a
multiplication by $i$.  The retained bound of four coefficient operations
per child therefore applies.  Its computation of address parities is
descriptor work.  The complete nonrecursive node allowance is
\[
 D_{\rm node}=D_{\rm scalar}+4s<64(W+m+1)^3,                    \tag{PS11b}
\]
where the strict inequality follows by integer comparison at the values
below.  Source, sink and bank corrections were already included in the
$8W$ term; the original bound for them is $4W+4\le8W$.

For completeness, the number of edges also stays explicitly bounded.
Refine each grouped scalar gate into two-role shears, retaining its common
label throughout that refinement.  Its number of touched-role incidences
is at most twice the operation count used above.  An identity incidence
may be retained for each side role in each of the four mixer applications.
A physical role participates in at most two shared invocations, so there
are at most $8W$ such additional incidences.  Each role contributes one
more edge than its number of gate incidences.  Hence the total edge count
is at most
\[
 Q_{\rm edge}=2D_{\rm scalar}+9W<100W.                          \tag{PS11c}
\]
All new edges created by a refinement have zero residual.  For every
nonzero residual edge its two binary changes of basis need at most
$2(m(m-1)+3m)$ binary row additions.  These operations permute complete
coefficient encodings and add zero coefficient depth.  They are charged
to the compact-control movement bound in time, not to magnitude or
denominator growth.  Thus they require no additional term in (PS11b).

Consequently the retained guard recurrence may use its existing node
charge $E=64(W+m+1)^3$, now evaluated with the new network constants.
The recursive residual calls are charged separately by $sA(e/m)$,
as before.  Address permutations and compact-control repair introduce
no additional coefficient arithmetic.

\paragraph{Finite counts and the phase interface.}
There are $3v^2$ invocations.  First/third-stage reuse leaves two complete
auxiliary banks per pair of fixed triples, giving
\[
 W=2N+2v^2(R_{\rm side}+h),\qquad
 L_{\rm dec}=3v^2h^2,\qquad
 s=Wm-2N+2L_{\rm dec}.                                         \tag{PS11}
\]
Here $L_{\rm dec}$ still counts central decreases in all three stages;
sharing a central role does not remove either of its two separate returns.
The signed sum of edge dimension changes telescopes to $Wm-2N$, and
every decrease has just been listed, proving the formula for $s$.

At $h=24$, exhaustive exact support generation by the stated mask rule gives
\[
\begin{gathered}
 v=2024,\quad m=13824,\quad N=8291469824,\\
 c_{\rm add}=77454,\quad q_{\rm out}=469294,\quad
 R_{\rm side}=546748,\\
 W=4496369044992,\quad L_{\rm dec}=7078883328,\\
 s=62157803252796416,\qquad Wm-s=2425172992>0.
\end{gathered}                                                  \tag{PS12}
\]
The companion checker verifies every mask split, every designated signed
output, the complete side matrix, and these integer counts.  These finite
checks support the explicit graph and its arithmetic; the general frame
and tape claims use the arguments given here and the retained interfaces.

\begin{proposition}\label{prop:pair-star-complex-interface}
The $h=24$ construction has $W$ complex roles and total binary phase-residual
length $s$ given in (PS12).  It performs the original signed bank exchange,
restores arbitrary auxiliary values, and satisfies the original complex
phase interface on any positive number of columns.  In particular
\[
 \frac{s}{W}<m^\sigma,\qquad \sigma=1-\frac4{10^9}.
\]
\end{proposition}
\begin{proof}
The signed scalar schedule was proved in (PS4)--(PS5).  Every gate has a
common nondegenerate frame, every edge joins comparable labels, and every
nonzero residual has an orthonormal basis by (PS8)--(PS10).  The original
binary phase factorization therefore contributes exactly one forward or
inverse translation kernel per residual basis vector, totaling $s$.
The input/output data labels remain $U_a,0$ and
$\mathcal F,U_a^\perp$ for $X_a,Y_a$, respectively.  Every surviving
auxiliary role starts at zero, ends at $\mathcal F$, and is fixed by the
scalar routing.  The original signed endpoint correction consequently
applies on every role.  More explicitly, $u_a$ has norm one and weight
$27$; the $X_a\to Y_a$ route has the same 27 inverse factors per column.
The original source and sink $Z$ corrections, including the scalar
$i^{27k}$ on $k$ columns, turn these into forward factors.  Negating the
physical $X$ outputs and undoing the bank permutation then gives the
required $C^{\otimes mk}$ on every role.  Finally the exact arithmetic
comparisons
\[
 \sum_{j=0}^{15}\frac{(48/5)^j}{j!}>13824=m,\qquad
 \eta:=\frac{Wm-s}{Wm}=\frac{37}{948319488}
                   >\frac4{10^9}\frac{48}{5}
\]
give $\log m<48/5$ and $\eta>(4/10^9)\log m$.
The strict inequality $e^{-x}>1-x$ for $x>0$ now gives
$m^{1-4/10^9}>m(1-\eta)=s/W$.
\end{proof}
